\section{Regularidad espacial}\label{sec:regularidad}

Fijemos ahora la masa y desplacémonos de una línea a otra de la misma familia.
Al trasladar una diagonal cambia también el segmento que queda dentro del cuadrado.
Por eso la estimación debe controlar a la vez la variación de los factores y
el efecto de los extremos del segmento.

La regularidad se obtiene comparando el perfil con sus trasladados mediante la coercividad de una
formulación dual, sin suponer derivadas del maximizador. En $\cH=\bigoplus_fL^2(I_f)$ usamos la
norma $\|w\|_\cH^2=\sum_f\|w_f\|_{L^2(I_f)}^2$.

\subsection{Actividad fija}

Fijamos $\rho\in(0,1)$, $\lambda=\lambda(\rho)$, $\nu=\nu_\rho$ y las versiones continuas del
Lema~\ref{lem:vacancias}: $u_f=1/(1+\lambda J_f(u))$, con $J_f(u)$ la integral sobre la fibra de
$\prod_{g\ne f}u_g\circ\ell_g$, $0\le J_f\le1$ y $\kappa_\lambda:=(1+\lambda)^{-1}\le u_f\le1$. Para
campos acotados $v_f=\log u_f$ definimos
\begin{equation}\label{eq:dual}
 \cE_\lambda(v)=\sum_f\int_{I_f}(e^{v_f}-v_f)+\lambda\int_{[0,1]^2}\exp\Bigl(\sum_fv_f\circ\ell_f\Bigr).
\end{equation}
Su gradiente en $\cH$ es $\cG_{\lambda,f}(v)=e^{v_f}-1+T_f\nu_v$, con
$\nu_v=\lambda\exp(\sum_fv_f\circ\ell_f)$, y el campo del maximizador cumple $\cG_\lambda(v)=0$; no se
impone masa a los campos de comparación, sino la actividad. Las derivaciones se justifican por
acotación y convergencia dominada a lo largo de segmentos de campos acotados. Para una dirección
acotada $w$,
\begin{equation}\label{eq:hess}
 \langle w,D\cG_\lambda(v)w\rangle=\sum_f\int e^{v_f}w_f^2+\int\nu_v\Bigl(\sum_fw_f\circ\ell_f\Bigr)^2 .
\end{equation}

\begin{lema}[Estabilidad]\label{lem:estabilidad}
Sea $0<\kappa'\le1$. Si $u=e^v$ y $\widetilde u=e^{\widetilde v}$ toman valores en $[\kappa',1]$,
entonces $\langle\cG_\lambda(\widetilde v)-\cG_\lambda(v),\widetilde v-v\rangle\ge\kappa'\|\widetilde v-v\|_\cH^2$.
En particular, si $\cG_\lambda(v)=0$, entonces
$\|\widetilde v-v\|_\cH\le\|\cG_\lambda(\widetilde v)\|_\cH/\kappa'$.
\end{lema}

\begin{proof}
El segmento entre los logaritmos tiene exponenciales en $[\kappa',1]$. Se integra~\eqref{eq:hess} a lo
largo de él y se aplica Cauchy--Schwarz.
\end{proof}

Extendemos $u_D$ y $u_A$ por uno fuera de $I_D$ e $I_A$ (el empalme es continuo por el
Lema~\ref{lem:vacancias}) y escribimos $\widehat u_D,\widehat u_A$. Para $t\in\mathbb R$ sea
\[
 \widehat u_R(t)=\Bigl[1+\lambda\int_0^1u_C(c)\,\widehat u_D(t-c)\,\widehat u_A(t+c)\,dc\Bigr]^{-1},
\]
que coincide con $u_R$ en $[0,1]$ y toma valores en $[\kappa_\lambda,1]$. Para un desplazamiento $t$,
el campo comparado, en los dominios originales, es
\[
 u_R^t(r)=\widehat u_R(r+t),\quad u_C^t=u_C,\quad u_D^t(d)=\widehat u_D(d+t),\quad
 u_A^t(s)=\widehat u_A(s+t),\qquad v^t=\log u^t .
\]

\begin{lema}[Residuo de la traslación]\label{lem:residuo}
Con $\eta^t=\cG_\lambda(v^t)$ se tiene $\eta^t_R=0$, $|\eta^t_C|,|\eta^t_D|,|\eta^t_A|\le2\lambda|t|$ y
$\|\eta^t\|_\cH\le2\sqrt5\,\lambda|t|$. En consecuencia $\|v^t-v\|_\cH\le2\sqrt5\,\lambda(1+\lambda)|t|$.
\end{lema}

\begin{proof}
La primera igualdad es la definición de $\widehat u_R$. En una columna $c$, se sustituye $r'=r+t$ en
la integral de $u_R^t(r)\,u_D^t(r-c)\,u_A^t(r+c)$: los integrandos coinciden con los originales donde
ambos dominios coinciden, y $r'\in[0,1]$ se cambia por $r'\in[t,1+t]$, cuya diferencia simétrica mide
a lo sumo $2|t|$; los productos de vacancias tienen módulo a lo sumo uno. En la diagonal $d$, tras la
misma sustitución, $c=r'-(d+t)$ y el integrando es $\widehat u_R(r')\,u_C(c)\,\widehat u_A(2r'-(d+t))$;
el dominio trasladado exige $r'\in[t,1+t]$ y $c\in[0,1]$, y el original de la diagonal $d+t$ exige
$r'\in[0,1]$ y $c\in[0,1]$, de modo que su diferencia está en las mismas dos tiras. Si $d+t\notin I_D$,
la fibra original es vacía y $\widehat u_D(d+t)=1$, y el argumento sigue valiendo. Para $A$ se usa
$c=s+t-r'$. El residuo es $\lambda u^t_f$ por la diferencia de las dos integrales, y su módulo es a lo
sumo $2\lambda|t|$. Los dominios de $C$, $D$ y $A$ miden $1$, $2$ y $2$, lo que da la cota en $\cH$; la
última desigualdad es el Lema~\ref{lem:estabilidad} con $\kappa'=\kappa_\lambda$.
\end{proof}

Para la traslación en la segunda coordenada se intercambian $R$ y $C$: se define $\widehat u_C$ por su
ecuación y se comparan $u_R(r)$, $\widehat u_C(c+t)$, $\widehat u_D(d-t)$ y $\widehat u_A(s+t)$; la
componente $C$ del residuo es nula y las otras tres cumplen las mismas cotas.

Para pasar de $L^2$ a diferencias puntuales: si $u=e^v$ y $\widetilde u=e^{\widetilde v}$ toman valores en
$[0,1]$, telescopar el producto de tres factores en una fibra da
\begin{equation}\label{eq:conversion}
 |J_f(\widetilde u)-J_f(u)|\le\sum_{g\ne f}\|\widetilde u_g-u_g\|_{L^2(I_g)}\le\sqrt3\,\|\widetilde v-v\|_\cH ,
\end{equation}
porque en una fibra de $f$ el parámetro $\ell_g$, $g\ne f$, avanza con pendiente de módulo uno o dos
(nunca cero); Cauchy--Schwarz y la longitud paramétrica $\le1$ dan la primera desigualdad, y
$|e^\alpha-e^\beta|\le|\alpha-\beta|$ para $\alpha,\beta\le0$ la segunda. Para el campo trasladado,
$u^t_f(1+\lambda J_f(u^t))=1+\eta^t_f$; comparando con $u_f(1+\lambda J_f(u))=1$, con denominadores
$\ge1$, el Lema~\ref{lem:residuo} y~\eqref{eq:conversion},
\begin{equation}\label{eq:traslacion}
 |u^t_f-u_f|\le2\lambda|t|+2\sqrt{15}\,\lambda^2(1+\lambda)|t|\le L(\lambda)|t|,\qquad
 L(\lambda)=2\lambda+8\lambda^2(1+\lambda).
\end{equation}
La traslación de filas controla $R$, $D$ y $A$, y la de columnas controla $C$.
